\documentclass[10pt]{article}
\usepackage[a4paper,margin=0.72in]{geometry}
\usepackage{array}
\usepackage{amsmath}
\usepackage{booktabs}
\usepackage{enumitem}
\usepackage{hyperref}
\usepackage[table]{xcolor}
\usepackage{titlesec}

\hypersetup{colorlinks=true,linkcolor=black,urlcolor=blue,citecolor=black}
\setlength{\parindent}{0pt}
\setlength{\parskip}{3pt}
\setlist[itemize]{leftmargin=1.1em,itemsep=1pt,topsep=2pt}
\setlist[enumerate]{leftmargin=1.25em,itemsep=1pt,topsep=2pt}
\renewcommand{\arraystretch}{1.13}
\titleformat{\section}{\large\bfseries}{\thesection}{0.45em}{}
\titlespacing*{\section}{0pt}{6pt}{3pt}
\titleformat{\subsection}{\normalsize\bfseries}{\thesubsection}{0.45em}{}
\titlespacing*{\subsection}{0pt}{5pt}{2pt}

\newcommand{\blank}[1][8em]{\colorbox{yellow!35}{\rule{0pt}{1.1em}\hspace{#1}}}
\newcommand{\eg}{E_{\mathrm{g}}}
\newcommand{\ehull}{E_{\mathrm{hull}}}
\newcommand{\form}{x_{\mathrm{form}}}
\newcommand{\struc}{x_{\mathrm{struc}}}
\newcommand{\unc}{u}
\newcommand{\domain}{d_{\mathrm{dom}}}
\newcommand{\utility}{U}

\begin{document}

\begin{center}
{\Large \textbf{AI-Driven Design of Functional Boron Nitride Materials}}\\
{\normalsize Research Plan}
\end{center}

\textbf{Principal investigator, position and affiliated department}\\
Prof. Ma Chen, Assistant Professor, Department of Computer Science, City University of Hong Kong.

\textbf{Collaborator(s) and respective department(s), if any}\\
Shenzhen Zhongneng Tianyi Technology Co., Ltd. (industrial collaborator).

\textbf{Duration of the project}\\
\blank[18em]

\section*{Abstract}
Boron nitride (BN) materials are promising for wide-band-gap ultraviolet optoelectronics and dielectric support layers in two-dimensional heterostructures. This project aims to build an uncertainty-aware AI pipeline to prioritize BN-related materials for these two application targets. The study will construct a bounded BN design space, learn property and uncertainty models from public materials databases, rank candidates under band-gap, stability, novelty, and domain-distance constraints, and generate first-pass structure hypotheses for selected candidates. Expected outcomes include a curated BN-centered dataset, benchmarked property models, an uncertainty-aware ranked candidate table, structure handoff files, and manuscript-ready analysis for academic and industrial collaboration.

\section*{Key objectives}
\begin{enumerate}
  \item \textbf{Define a bounded BN design space.}
  \[
    \mathcal{C}_{\mathrm{BN}}=
    \mathcal{C}_{\mathrm{core}}\cup\mathcal{C}_{\mathrm{ext}}\cup
    \mathcal{C}_{\mathrm{explore}},
  \]
  where $\mathcal{C}_{\mathrm{core}}=\{\mathrm{BN},\mathrm{h\mbox{-}BN},
  \mathrm{c\mbox{-}BN},\mathrm{doped\ BN}\}$,
  $\mathcal{C}_{\mathrm{ext}}=\{\mathrm{BCN},\mathrm{BC_2N},\mathrm{SiBN}\}$,
  and $\mathcal{C}_{\mathrm{explore}}$ contains same-group substituted analogues
  limited to chemically plausible substitutions around B, N, C, and Si sites.
  \textit{Output:} a curated dataset with formula, structure, target-property,
  and provenance fields.

  \item \textbf{Learn property and uncertainty models.}
  \[
    f_{\theta}:(\form,\struc)\rightarrow
    \mathbf{y}=(\eg,\Delta H,\ehull,\epsilon_r,\rho_{\mathrm{2D}},\unc).
  \]
  \textit{Output:} benchmarked models for band gap, stability proxy, dielectric
  relevance, and predictive uncertainty.

  \item \textbf{Prioritize candidates under application constraints.}
  \[
    \utility(c)=w_1\hat{\eg}(c)-w_2\hat{\ehull}(c)-w_3\hat{\unc}(c)
    -w_4\domain(c)+w_5\rho_{\mathrm{novel}}(c)+w_6\rho_{\mathrm{app}}(c).
  \]
  \textit{Output:} a ranked table of candidate families with predicted
  properties, uncertainty, novelty score, and action labels.

  \item \textbf{Generate structure hypotheses for selected candidates.}
  \[
    c \mapsto \{s_{c,1},\ldots,s_{c,k}\}
    \mapsto \{\mathrm{geometry},\mathrm{relaxation},\mathrm{property\;recheck}\}.
  \]
  \textit{Output:} structure files and a validation handoff package for
  high-priority candidates.
\end{enumerate}

\section*{Research methodology}
\subsection*{WP1. Data construction and cleaning}
The project starts from public computational materials resources rather than an unbounded chemical search. 2DMatPedia provides more than 6,000 monolayer structures generated by top-down and bottom-up strategies \textsuperscript{[1]}. JARVIS-Tools supports programmatic access to atomistic materials datasets \textsuperscript{[2]}, while Matbench, matminer, the Materials Project, and C2DB provide benchmark, featurization, and reference spaces for materials ML \textsuperscript{[3--6]}. Records will be deduplicated, aligned to target-property definitions, and labelled with provenance:

\[
  r_i=(c_i,s_i,\mathbf{y}_i,p_i),\qquad
  \mathcal{D}_{\mathrm{BN}}=\{r_i:c_i\in\mathcal{C}_{\mathrm{BN}}\},\qquad
  p_i=(\mathrm{source},\mathrm{method},\mathrm{version}).
\]
This creates the BN-centered data layer that supports all later modelling and ranking decisions.

\subsection*{WP2. Property prediction and uncertainty estimation}
The modelling target is a decision vector rather than a single benchmark number:
\[
  \max_{c,s}\; \mathbf{g}(c,s)=
  \big[\eg,\;-\ehull,\;\rho_{\mathrm{dielectric}},\;\rho_{\mathrm{UV}},
  \rho_{\mathrm{2D}},\;-\unc\big].
\]
Model families will include descriptor baselines, composition neural networks, and crystal graph models, drawing on established materials-ML tools such as pymatgen, CGCNN, MEGNet, ALIGNN, CrabNet, and CHGNet \textsuperscript{[7--12]}. Predictive uncertainty will be estimated using model ensembles and calibrated against validation-set errors:
\[
  \hat{\unc}_{\mathrm{cal}}(c)=
  \alpha\cdot \mathrm{Std}_{m\in\mathcal{M}}\big[f_m(c)\big]+
  \beta\cdot \widehat{\mathrm{err}}_{\mathrm{val}}(c).
\]
Candidates with high predicted performance but high uncertainty will be labelled as exploratory rather than priority candidates. These models are not the novelty by themselves; they provide a controlled basis for a BN-specific decision pipeline.

\subsection*{WP3. Multi-objective ranking and candidate selection}
The output for each candidate will be a decision vector
\[
  \mathbf{z}(c)=
  (\hat{\eg},\hat{\ehull},\hat{\epsilon}_r,\hat{\unc},\domain,
  \rho_{\mathrm{novel}},\rho_{\mathrm{app}},a),
\]
where domain distance will be measured using composition descriptors, structural descriptors, and learned embedding distances from the BN-centered training set:
\[
  \domain(c)=\min_{r_i\in\mathcal{D}_{\mathrm{BN}}}
  \left\|\phi(c)-\phi(c_i,s_i)\right\|_{\Sigma}.
\]
Novelty will be estimated by distance from known BN-family compounds, while application relevance will be assigned according to whether the candidate matches the wide-gap/UV or dielectric/2D-support target. The action label is
\[
  a\in\{\mathrm{control},\mathrm{priority},\mathrm{explore},\mathrm{hold}\}.
\]
Candidates will be prioritized only when predicted performance, uncertainty, and domain-distance criteria are jointly satisfied:
\[
  a(c)=\mathrm{priority}\quad\mathrm{if}\quad
  \utility(c)\ge q_{\alpha},\quad \hat{\unc}(c)\le u_{\max},\quad
  \domain(c)\le d_{\max}.
\]
This connects prediction to research action: reference controls, high-priority validation, exploratory families, and candidates reserved for later evidence.

The candidate selection protocol is: generate a BN-related candidate pool; predict target properties; estimate uncertainty and domain distance; filter by stability, uncertainty, and application target; assign action labels.

\subsection*{WP4. Structure generation, validation handoff, and iteration}
Selected formulas will first be converted into candidate structures through prototype reuse, database-derived structure transfer, and controlled substitution. Generative models such as CDVAE will be used only when prototype-based routes are insufficient \textsuperscript{[13]}. A first-pass handoff package will include structure files, predicted properties, uncertainty scores, provenance, and recommended validation actions.

The validation protocol will test both interpolation and extrapolation. Models will be evaluated using random splits, grouped-by-formula splits, and BN-family-aware splits. Candidate ranking will be assessed by model ensembles and perturbation of objective weights:
\[
  \mathrm{MAE},\ \mathrm{RMSE},\ R^2,\ \tau_{\mathrm{Kendall}},
  \rho_{\mathrm{Spearman}},\ P(c\in\mathrm{Top}\text{-}k).
\]
For selected priority candidates, materials-side validation will include charge-neutrality checks, structural sanity checks, geometry relaxation where feasible, energy-above-hull re-estimation,
\[
  \ehull(c,s)=E(c,s)-E_{\mathrm{hull}}\big(\mathrm{composition}(c)\big),
\]
and comparison with known BN-family references.
The current local project already provides a starting implementation for public-data ingestion, BN-themed candidate generation, model comparison, rank stability, and first-pass structure handoff. The proposed research will extend this base into a larger, uncertainty-aware BN prioritization platform.

\section*{Significance and major deliverables}
BN is a high-value family for wide-band-gap ultraviolet materials, dielectric layers, and 2D heterostructure platforms \textsuperscript{[14--16]}. Thermal management remains a secondary relevance, while the main proposal focuses on wide-gap/UV and dielectric/2D-support targets. The scientific question is:
\[
  \textit{Which BN-related compositions and structures should be investigated next, and with what level of uncertainty?}
\]
The novelty lies in the BN-centered data layer, uncertainty-aware candidate selection, multi-objective ranking under application constraints, formula-to-structure handoff workflow, and rank-stability analysis for materials prioritization.

\textbf{Major deliverables}
\begin{enumerate}
  \item \textbf{D1:} BN-centered dataset with formula, structure, target-property, and provenance fields.
  \item \textbf{D2:} Benchmark report comparing descriptor, composition, and graph-based models.
  \item \textbf{D3:} Ranked candidate table including formula, predicted band gap, stability proxy, uncertainty, domain distance, novelty score, application label, and recommended action.
  \item \textbf{D4:} Structure hypothesis package including CIF/POSCAR files, provenance record, predicted properties, uncertainty score, and validation notes.
  \item \textbf{D5:} Technical report, figures, demo, and manuscript-ready analysis.
\end{enumerate}

\section*{Ethics and safety}
The initial project is computational and data-driven. It uses public databases, software, models, and generated candidate structures. If later work involves chemical synthesis or regulated laboratory procedures, the responsible experimental party will handle the required safety approvals.

\section*{Selected references}
\small
\begin{enumerate}[leftmargin=1.2em,itemsep=1pt]
  \item Zhou, J. et al. 2DMatPedia: an open computational database of two-dimensional materials from top-down and bottom-up approaches. \textit{Scientific Data} 6, 86 (2019).
  \item Choudhary, K. et al. JARVIS-Tools: open-access software for atomistic data-driven materials design. \textit{npj Computational Materials} 6, 173 (2020).
  \item Dunn, A. et al. Benchmarking materials property prediction methods: the Matbench test set and Automatminer reference algorithm. \textit{npj Computational Materials} 6, 138 (2020).
  \item Ward, L. et al. Matminer: An open source toolkit for materials data mining. \textit{Computational Materials Science} 152, 60--69 (2018).
  \item Jain, A. et al. Commentary: The Materials Project: a materials genome approach to accelerating materials innovation. \textit{APL Materials} 1, 011002 (2013).
  \item Haastrup, S. et al. The Computational 2D Materials Database: high-throughput modeling and discovery of atomically thin crystals. \textit{2D Materials} 5, 042002 (2018).
  \item Ong, S. P. et al. Python Materials Genomics (pymatgen): A robust, open-source python library for materials analysis. \textit{Computational Materials Science} 68, 314--319 (2013).
  \item Xie, T. and Grossman, J. C. Crystal graph convolutional neural networks for an accurate and interpretable prediction of material properties. \textit{Physical Review Letters} 120, 145301 (2018).
  \item Chen, C. et al. Graph networks as a universal machine learning framework for molecules and crystals. \textit{Chemistry of Materials} 31, 3564--3572 (2019).
  \item Choudhary, K. and DeCost, B. Atomistic line graph neural network for improved materials property predictions. \textit{npj Computational Materials} 7, 185 (2021).
  \item Wang, A. Y.-T. et al. CrabNet for explainable deep learning in materials science. \textit{npj Computational Materials} 7, 77 (2021).
  \item Deng, B. et al. CHGNet as a pretrained universal neural network potential for charge-informed atomistic modelling. \textit{Nature Machine Intelligence} 5, 1031--1041 (2023).
  \item Xie, T. et al. Crystal diffusion variational autoencoder for periodic material generation. \textit{ICLR} (2022).
  \item Watanabe, K., Taniguchi, T. and Kanda, H. Direct-bandgap properties and evidence for ultraviolet lasing of hexagonal boron nitride single crystal. \textit{Nature Materials} 3, 404--409 (2004).
  \item Cassabois, G., Valvin, P. and Gil, B. Hexagonal boron nitride is an indirect bandgap semiconductor. \textit{Nature Photonics} 10, 262--266 (2016).
  \item Dean, C. R. et al. Boron nitride substrates for high-quality graphene electronics. \textit{Nature Nanotechnology} 5, 722--726 (2010).
\end{enumerate}
\normalsize

\section*{Project Details}
\textbf{Project Title:} AI-Driven Design of Functional Boron Nitride Materials

\textbf{PI (Dept):} Prof. Ma Chen (Department of Computer Science)

\textbf{Budget Breakdown}
\begin{center}
\small
\begin{tabular}{p{0.20\linewidth}p{0.51\linewidth}p{0.17\linewidth}}
\toprule
\textbf{Budget Item} & \textbf{Detail breakdown} & \textbf{Amount (HKD)} \\
\midrule
Staffing & Student helper / research assistant for data curation, experiments, documentation, and demo maintenance. & \blank \\
Equipment & Compute, storage, and structure-validation resources. & \blank \\
General expenses & Cloud credits, backup, software services, hosting, and data-management costs. & \blank \\
Conference & Presentation at AI-for-science, materials informatics, or applied computing venues. & \blank \\
Overseas Travel & Optional collaborator visit, conference travel, or partner technical meeting. & \blank \\
\midrule
 & \textbf{Total Funding} & \blank \\
\bottomrule
\end{tabular}
\end{center}

\end{document}
